%  Esp_Loader: the board flashes ANOTHER ESP32 -- the esptool serial protocol
%  as a device-side client, with the auto-reset circuit done in software and
%  the result proven by the target's own hash.
\chapter[Esp\_Loader: Flashing Another ESP32]%
        {Esp\_Loader: Flashing Another ESP32}
\label{ch:esp_loader}

Every ESP32 ships with a ROM bootloader that accepts firmware over a serial
port. That is what \texttt{esptool} talks to from a PC. \texttt{ESP32S3.Esp\_Loader}
is the same protocol as a device-side client, so a board built on this SDK can
program \emph{other} ESP32s itself --- no PC in the
loop.\index{Esp\_Loader}\index{ROM bootloader}\index{esptool}

The shape of the job: reset the target into its ROM loader, synchronise,
describe the flash, stream the image in 1\,KB blocks, and reset the target into
its new firmware. Each step is a command with a reply, framed in SLIP. The
chapter walks the stack from the wire up.

\section{The wire: SLIP frames and commands}\index{SLIP framing}

The ROM speaks a request/response protocol over 115\,200-baud serial. Every
packet travels inside a SLIP frame: a \texttt{C0} delimiter on each end, with
\texttt{C0} and \texttt{DB} bytes inside the payload escaped. A request starts
with a fixed header --- direction (0), opcode, payload length, checksum --- and
the reply echoes the opcode back with a status trailer.

Two details vary by chip, and both change how bytes are \emph{parsed}, not just
what they mean. The original ESP32's ROM ends replies with four status bytes;
everything since uses two. And chips after the ESP32 take a fifth word in
\texttt{FLASH\_BEGIN}. The driver therefore identifies the target first and
keeps both facts in the session.

\section{Reset into the loader: the circuit esptool expects}
\index{auto-reset}\index{DTR/RTS}

A target enters its ROM loader when it leaves reset with \texttt{IO0} held low.
Development boards do this with two transistors cross-coupled between the serial
port's DTR and RTS lines, plus an RC on the EN pin. \texttt{esptool} wiggles
DTR and RTS in a fixed dance and trusts that circuit to translate.

\texttt{Esp\_Loader.Auto\_Reset} is that circuit in software, for a board that
sits between a PC and a target as a USB-serial bridge. The cross-coupling
matters: terminal programs assert both lines when they open a port, and the
real circuit's both-asserted state deliberately does nothing --- otherwise every
\texttt{screen} session would reset the target. The RC matters too.
\texttt{esptool}'s classic reset passes through a brief both-asserted moment,
and only the capacitor's hold-up carries EN across it. The software circuit
reproduces both: asserting reset is immediate, releasing it waits out an
emulated RC, and a release still pending when reset is wanted again is simply
cancelled.

When the board itself is the programmer, none of that subtlety is needed:
\texttt{Reset\_Target} drives the two lines directly, holds reset for
100\,ms with boot asserted, and releases them in order.

\section{The transport: a Link, not a UART}\index{Link record}

The protocol core does not name a UART. It talks through a \texttt{Link} --- a
record of access-to-subprogram values plus an opaque context, the same pattern
as \texttt{Block\_Dev}:

\begin{lstlisting}
type Link is record
   Ctx          : System.Address := System.Null_Address;
   Send         : Send_Proc    := null;
   Receive      : Receive_Proc := null;
   Assert_Reset : Line_Proc    := null;   --  logical, not electrical
   Assert_Boot  : Line_Proc    := null;
   Set_Baud     : Baud_Proc    := null;   --  optional
end record;
\end{lstlisting}

\texttt{Serial\_Link} is the ready-made transport: a UART plus two GPIOs, with
the electrical polarity of the reset and boot lines a configuration fact
(\texttt{Reset\_Drives\_High}), because a board that drives a transistor
inverts the sense. Its \texttt{Open} holds the target in reset while the UART
comes up, by default. That ordering is not politeness --- a UART FIFO reset
that races live traffic can leave the receiver serving rotated data for the
rest of the boot (Chapter~\ref{ch:hal} tells that story). A caller that must
not disturb a running target opts out and accepts the hazard knowingly.

\section{A session, start to finish}\index{FLASH\_BEGIN}\index{SPI\_ATTACH}

\begin{lstlisting}
Open (Wire, On => UART1, Tx => 4, Rx => 5,
      Reset => 6, Boot => 18);
Connect (Session, As_Link (Wire), Status);   --  reset, sync, identify
Configure_Flash (Session, 4 * 1024 * 1024, Status);
Begin_Image (Session, At_Offset => 16#1_0000#,
             Length => Size, Status => Status);
loop
   ... read a chunk from wherever the image lives ...
   Write (Session, Chunk, Status);
end loop;
End_Image (Session, Status);
Finish (Session, Run => True, Status => Status);  --  boot it
\end{lstlisting}

\texttt{Connect} retries the whole reset-and-sync sequence a few times, because
a ROM still waking up misses the first sync. It then drains the line until the
target goes quiet: the ROM answers one sync many times over, and a command sent
while those replies are still arriving reads against a stream that is out of
step. \texttt{Configure\_Flash} attaches the target's SPI flash and declares
its size (\texttt{SPI\_ATTACH}, then \texttt{SPI\_SET\_PARAMS}).
\texttt{Begin\_Image} erases the region --- the slowest step of a run --- and
\texttt{Write} streams it in the ROM's fixed 1\,KB blocks, re-blocking whatever
chunk sizes the caller finds convenient.

Identification deserves a note. \texttt{Connect} asks the modern way first
(\texttt{GET\_SECURITY\_INFO} carries a chip id on everything from the
ESP32-S3 on), and falls back to a magic register that tells the older chips
apart. A chip newer than the tables still connects, as \texttt{Unknown} ---
and \texttt{Unknown} deliberately selects the modern protocol variants, so an
unrecognised future chip is more likely to program than not. A caller that
only wants to write an image can skip identification entirely
(\texttt{Identify => False}); the reference host flasher never identifies at
all.

\section{Proving the write: the target hashes itself}
\index{SPI\_FLASH\_MD5}\index{MD5}\index{verification}

Every \texttt{FLASH\_DATA} block is checksummed and acknowledged, but a block
ack only says the target \emph{accepted} the bytes. The ROM's
\texttt{SPI\_FLASH\_MD5} command closes the loop: the target hashes a region of
its own flash and returns the digest. \texttt{Read\_Md5} exposes it, and
\texttt{ESP32S3.MD5} computes the expected side --- a streaming RFC~1321
implementation, pure Ada, no heap, added to the HAL for exactly this pairing.
Hash the image as it streams out, ask the target to hash what its flash holds,
and compare:

\begin{lstlisting}
MD5.Update (Digest, Chunk);        --  ... during the write loop
...
Read_Md5 (Session, Offset, Length, Readback, Status);
if Readback /= MD5.Hex_Digest (Digest) then ...  --  not programmed
\end{lstlisting}

MD5 is long broken as a cryptographic hash; this is an integrity check against
accident, not an adversary, and the ROM chose the algorithm. One sequencing
rule, learned the measured way: write every region first, then verify every
region. The ROM refuses a \texttt{FLASH\_BEGIN} that follows an MD5 command,
so the two passes must not interleave --- which is also the order
\texttt{esptool} uses.

\section{Testing without a victim}\index{fake ROM}\index{host testing}

The protocol core is host-testable. \texttt{test/esp\_loader\_host} runs the
real Ada client against a Python fake ROM over pipes: sync, identification for
each chip in the tables, byte-exact flashing, the auto-reset sequences as
\texttt{esptool} performs them, and the failure paths --- a silent target, a
refused command, an image shorter than declared. The fake ROM models each
chip's real reply shape, and one lesson from the bring-up is recorded in its
history: model what a trace of real silicon shows, not what documentation
implies. A DevKit and \texttt{esptool --trace} settle such questions in
minutes.

The harness proves the protocol; it cannot prove the transport. The one real
bug of the bring-up lived below every layer the fake ROM exercises --- an RX
FIFO with skewed pointers serving rotated bytes (Chapter~\ref{ch:hal}) --- and
survived precisely because the protocol logic above it was correct.

\medskip
\noindent One application built on this SDK closes the chapter's loop: a
programmer board that carries images on its own storage, flashes a blank
ESP32-S3 through \texttt{Esp\_Loader}, and calls the job done only when the
target's own hashes match --- bootloader, partition table and application,
each region verified.
